% Part: turing-machines
% Chapter: undecidability
% Section: enumerating-tms

\documentclass[../../../include/open-logic-section]{subfiles}

\begin{document}

\olfileid{tur}{und}{enu}
\olsection{ટ્યુરિંગ મશીનોની પરિગણના}

\begin{explain}
બધાં ટ્યુરિંગ મશીનોનો ગણ !!{enumerable} છે એ આપણે બતાવી શકીએ.
દરેક ટ્યુરિંગ મશીનનું સાન્ત વર્ણન આપી શકાય છે, તેમાંથી આ
વાત ફલિત થાય છે. અવસ્થાઓનો ગણ અને ટેપની વર્ણમાળા સાન્ત
ગણો છે. સંક્રમણ વિધેય એ $Q \times \Sigma$થી
$Q \times \Sigma \times \{\TMleft, \TMright, \TMstay\}$માં
જતું આંશિક વિધેય છે; તેથી જે સાન્ત સંખ્યામાં નિવેશ-જોડીઓ માટે તે
વ્યાખ્યાયિત છે, તેમનાં મૂલ્યોની યાદી આપીને તેને પણ
નિર્દિષ્ટ કરી શકાય છે.

અત્યાર સુધીની વાત સાચી છે, છતાં અહીં એક સૂક્ષ્મ ભેદ છે.
ટ્યુરિંગ મશીનની વ્યાખ્યાએ અવસ્થાઓના ગણ અને ટેપની
વર્ણમાળામાં કેવા !!{element}s હોઈ શકે તેના પર કોઈ
પ્રતિબંધ મૂક્યો નથી. તેથી, દાખલા તરીકે, દરેક વાસ્તવિક
સંખ્યા માટે તકનીકી અર્થમાં એવું ટ્યુરિંગ મશીન છે જે તે
સંખ્યાને અવસ્થા તરીકે વાપરે છે. જોકે મશીનનું \emph{વર્તન}
કયા પદાર્થો અવસ્થાઓ અને ચિહ્નો તરીકે કામ કરે છે તેના પર
આધારિત નથી. \olref{fig:variants}નાં બે ટ્યુરિંગ મશીનો જુઓ.
\begin{figure}
\begin{center}
\begin{tikzpicture}[->,>=stealth',shorten >=1pt,auto,node distance=2.8cm,
                    semithick]
  \tikzstyle{every state}=[fill=none,draw=black,text=black]

  \node[initial,state]         (A)              {$q_0$};
  \node[state]         (B) [right of=A] {$q_1$};

  \path (A) edge [bend left] node {\TMtrans{\TMstroke}{\TMstroke}{\TMright}} (B)
        (B) edge [loop above] node {\TMtrans{\TMblank}{\TMblank}{\TMright}} (B)
            edge [bend left] node {\TMtrans{\TMstroke}{\TMstroke}{\TMright}} (A);
\end{tikzpicture}\\
\begin{tikzpicture}[->,>=stealth',shorten >=1pt,auto,node distance=2.8cm,
  semithick]
\tikzstyle{every state}=[fill=none,draw=black,text=black]

\node[initial,state]         (A)              {$s$};
\node[state]         (B) [right of=A] {$h$};

\path (A) edge [bend left] node {\TMtrans{A}{A}{\TMright}} (B)
(B) edge [loop above] node {\TMtrans{\TMblank}{\TMblank}{\TMright}} (B)
edge [bend left] node {\TMtrans{A}{A}{\TMright}} (A);
\end{tikzpicture}
\end{center}
\caption{\emph{સમ} મશીનનાં બે રૂપ}
\ollabel{fig:variants}
\end{figure}
આ બે આલેખ બે મશીનોને અનુરૂપ છે: $M$ની ટેપ-વર્ણમાળા
$\Sigma = \{\TMendtape,\TMblank,\TMstroke\}$ અને અવસ્થાઓનો
ગણ $\{q_0,q_1\}$ છે, જ્યારે $M'$ની વર્ણમાળા $\Sigma' =
\{\TMendtape,\TMblank,A\}$ અને અવસ્થાઓનો ગણ $\{s,h\}$
છે. પરંતુ તેમની બાકીની સૂચનાઓ સમાન છે: $n$ સંખ્યાના
$\TMstroke$ના ક્રમ પર $M$ ત્યારે અને માત્ર ત્યારે અટકે છે
જ્યારે $n$ સમ હોય; અને $n$ સંખ્યાના $A$ના ક્રમ પર $M'$
પણ ત્યારે અને માત્ર ત્યારે અટકે છે જ્યારે $n$ સમ હોય.
આપણે માત્ર $\TMstroke$નું નામ~$A$, $q_0$નું નામ~$s$ અને
$q_1$નું નામ~$h$ રાખ્યાં છે. આ દાખલાને સામાન્ય બનાવી
શકાય: ચિહ્નો અને અવસ્થાઓનાં આવાં નામ બદલવાથી એકમાંથી
બીજું મળે તો ટ્યુરિંગ મશીનોને સમાન માની શકીએ. હકીકતમાં,
આપણે મશીનનાં ચિહ્નો અને અવસ્થાઓને સીધાં ધન પૂર્ણાંકો
જ માની શકીએ: $\sigma_0$ને બદલે~$1$, $\sigma_1$ને
બદલે~$2$ વગેરે; $\TMendtape$ એ~$1$, $\TMblank$ એ~$2$
વગેરે. આ રીતે \emph{સમ} મશીન \olref{fig:standard-even}માં
દર્શાવેલું મશીન બને છે.
\begin{figure}
\[\begin{tikzpicture}[->,>=stealth',shorten >=1pt,auto,node distance=2.8cm,
  semithick]
\tikzstyle{every state}=[fill=none,draw=black,text=black]

\node[initial,state]         (A)              {$1$};
\node[state]         (B) [right of=A] {$2$};

\path (A) edge [bend left] node {\TMtrans{3}{3}{\TMright}} (B)
(B) edge [loop above] node {\TMtrans{2}{2}{\TMright}} (B)
edge [bend left] node {\TMtrans{3}{3}{\TMright}} (A);
\end{tikzpicture}
\]
\caption{માનક \emph{સમ} મશીન}
\ollabel{fig:standard-even}
\end{figure}
જે ટ્યુરિંગ મશીનની અવસ્થાઓ અને ચિહ્નો ધન પૂર્ણાંકો છે
તેને \emph{માનક} મશીન કહી શકીએ અને હવે પછી માત્ર માનક
મશીનો જ વિચારીએ.\footnote{``માનક મશીન'' પરિભાષા
સ્વયં માનક નથી.}

આપણે ટ્યુરિંગ મશીનોનો ગણ !!{enumerable} છે તે બતાવવું
હતું; ઉપરની ચર્ચા મુજબ માનક ટ્યુરિંગ મશીનોનો ગણ
!!{enumerable} છે તે બતાવવું પૂરતું છે. ધારો કે માનક
ટ્યુરિંગ મશીન $M = \tuple{Q, \Sigma, q_0, \delta}$
આપેલું છે. ધન પૂર્ણાંકોની સાન્ત શ્રેણી વડે તેનું વર્ણન
કેવી રીતે કરીએ? પહેલાં અવસ્થાઓની સંખ્યા, પછી અવસ્થાઓ,
ચિહ્નોની સંખ્યા, ચિહ્નો અને પ્રારંભિક અવસ્થાની યાદી આપીએ.
($M$ માનક હોવાથી આ બધાં ધન પૂર્ણાંકો છે એ યાદ રાખો.)
હવે~$\delta$નું શું? શક્ય દલીલોનો ગણ, એટલે કે
$\tuple{q,\sigma}$ જોડો, સાન્ત છે, કારણ કે $Q$ અને
~$\Sigma$ સાન્ત છે. તેથી~$\delta$માં રહેલી માહિતી એ
એવી બધી $5$ ઘટકોવાળી ક્રમયુક્ત શ્રેણીઓ
$\tuple{q, \sigma, q', \sigma', d}$ની
સાન્ત યાદી છે, જેમના માટે
$\delta(q,\sigma) = \tuple{q', \sigma', D}$; અહીં $d$
દિશા~$D$નો સંકેત આપતો અંક છે (જેમ કે~$\TMleft$ માટે
$1$,~$\TMright$ માટે $2$ અને~$\TMstay$ માટે $3$).

આ રીતે દરેક માનક ટ્યુરિંગ મશીનને ધન પૂર્ણાંકોની સાન્ત
યાદી વડે, એટલે કે શ્રેણી $s_M \in (\PosInt)^*$ વડે,
વર્ણવી શકાય છે. દાખલા તરીકે, માનક \emph{સમ} મશીનનો
સંકેત આ શ્રેણી છે:
\[
2, \underbrace{1, 2}_Q, 3, \overbrace{1, 2, 3}^\Sigma, 1, \underbrace{1, 3, 2, 3, 2}_{\delta(1,3) = \tuple{2,3,R}},
\overbrace{2, 2, 2, 2, 2}^{\delta(2,2) = \tuple{2,2,R}},
\underbrace{2, 3, 1, 3, 2}_{\delta(2,3) = \tuple{1,3,R}}.
\]
\end{explain}

\begin{thm}
$\Nat$થી~$\Nat$માં જતાં કેટલાંક વિધેયો ટ્યુરિંગ-સંગણનીય
નથી.
\end{thm}

\begin{proof}
ધન પૂર્ણાંકોની સાન્ત શ્રેણીઓનો ગણ~$(\PosInt)^*$
!!{enumerable} છે એ આપણે જાણીએ છીએ
(\cref{sfr:siz:zigzag:prob:posint-star}). તેથી માનક
ટ્યુરિંગ મશીનોનાં વર્ણનોનો ગણ, $(\PosInt)^*$નો ઉપગણ
હોવાથી, પોતે પણ પરિગણનીય છે. $\Nat$થી~$\Nat$માં જતું
દરેક ટ્યુરિંગ-સંગણનીય વિધેય કોઈ ટ્યુરિંગ મશીન વડે
(હકીકતમાં, અનેક મશીનો વડે) ગણાય છે. તેની અવસ્થાઓ અને
ચિહ્નોનાં નામ ધન પૂર્ણાંકો વડે બદલીને (ખાસ કરીને
$\TMendtape$ને~$1$, $\TMblank$ને~$2$ અને
$\TMstroke$ને~$3$ રાખીને) જોઈ શકાય છે કે દરેક
ટ્યુરિંગ-સંગણનીય વિધેય કોઈ માનક ટ્યુરિંગ મશીન વડે
ગણાય છે. એટલે $\Nat$થી~$\Nat$માં જતાં બધાં
ટ્યુરિંગ-સંગણનીય વિધેયોનો ગણ પણ પરિગણનીય છે.

બીજી તરફ, $\Nat$થી~$\Nat$માં જતાં બધાં વિધેયોનો
ગણ !!{enumerable} નથી (\cref{sfr:siz:red:prob:nat-nat}).
જો દરેક વિધેય કોઈ ટ્યુરિંગ મશીન વડે સંગણનીય હોત,
તો તેમનું સંગણન કરતાં મશીનોનાં બધાં વર્ણનોની યાદી
આપીને બધાં વિધેયોની પરિગણના કરી શકાત. તેથી કેટલાંક
વિધેયો ટ્યુરિંગ-સંગણનીય નથી.
\end{proof}

\begin{prob}
  અવસ્થાઓ અને વર્ણમાળાનાં ચિહ્નોની સ્પષ્ટ યાદી આપવી ન
  પડે એ રીતે ટ્યુરિંગ મશીનોનું વર્ણન કેવી રીતે આપી શકાય
  તે વિચારો. ``માનક'' મશીનની તમારી પોતાની વ્યાખ્યા આપી
  શકો છો; પરંતુ તમારા નવા અર્થમાં દરેક ટ્યુરિંગ મશીન
  જે વિધેય ગણે છે તે કોઈ ``માનક'' મશીન પણ કેમ ગણી
  શકે તે વિશે કંઈક કહો.
  \sourcecorrection{OLUN-004}{સ્થિર અંગ્રેજી મૂળમાં
  મશીનની જ સંગણના કરવાની વાત છે; અહીં અર્થ તેના
  અવસ્થા અને ચિહ્નોનાં નામ બદલીને સમાન સંગણન
  કરતું માનક મશીન મેળવવાનો છે.}
\end{prob}

\end{document}
